\section{Conclusion}

SC-ELAN provides a practical recipe for improving small object detection by jointly optimizing backbone representation and head-level class calibration.
Across 29 controlled runs on VisDrone2019-DET, we derive the following conclusions:
(i)~\textbf{selective long-range context} (LSKA with gating) consistently improves strict localization;
(ii)~\textbf{head reweighting is high-ROI but sensitive}, where the Soft CAI regime is stable while overly smoothed priors can regress;
and (iii)~\textbf{aggressive efficiency compression hurts tail classes first}.
Our best configuration (v2-P1d, $\alpha{=}0.10$, $\beta{=}0.40$) reaches \textbf{mAP50-95 = 0.212} on the VisDrone test-dev benchmark while maintaining real-time latency on RTX~4090.

\subsection{Limitations and Future Work}
We plan to (1) report multi-seed statistics for key variants to quantify robustness, (2) validate cross-dataset generalization on at least one additional aerial/traffic benchmark, and (3) further explore $\\beta>0.40$ under fixed $\\alpha=0.10$ to map the performance ceiling of CAI reweighting for long-tail categories.
